\documentclass[pdfa,a4paper,UKenglish,cleveref,autoref,thm-restate]{lipics-v2021}

\let\CheckCommand\providecommand

\bibliographystyle{plainurl}% the mandatory bibstyle

\title{Constructing (Co)inductive Types via Large Sizes}

\author{Bastiaan Laarakker}{LIACS, Leiden University, The Netherlands}{b.g.laarakker@liacs.leidenuniv.nl}{https://orcid.org/0009-0005-3079-8505}{This work was partially supported by the Dutch Research Council (NWO) under grant number\ VI.Veni.232.286 (ChEOpS).}
\author{Dani\"el Otten}{ILLC, University of Amsterdam, The Netherlands \and \url{http://otten.co}}{daniel@otten.co}{https://orcid.org/0000-0003-2557-3959}{}
\author{Benno van den Berg}{ILLC, University of Amsterdam, The Netherlands \and \url{https://staff.fnwi.uva.nl/b.vandenberg3/}}{b.vandenberg3@uva.nl}{}{}

\funding{\emph{Benno van den Berg and Dani\"el Otten:} This publication is part of the project “The Power of Equality” (with project
number OCENW.M20.380) of the research programme Open Competition Science M 2 which is financed by the Dutch Research Council (NWO).}

\hideLIPIcs
\nolinenumbers

\relatedversion{This paper is based on the master thesis of the first author, which was supervised by the second and third authors.}
% \relatedversiondetails{Full Version}{https://arxiv.org/abs/2602.18921}
\authorrunning{B. Laarakker, D. Otten and B. van den Berg } 
\Copyright{Bastiaan Laarakker, Benno van den Berg and Dani\"el Otten} 
\begin{CCSXML}
<ccs2012>
   <concept>
       <concept_id>10003752.10003790.10011740</concept_id>
       <concept_desc>Theory of computation~Type theory</concept_desc>
       <concept_significance>500</concept_significance>
       </concept>
 </ccs2012>
\end{CCSXML}

\acknowledgements{We are thankful to Andreas Abel, Jesper Cockx, and Andreas Nuyts for helpful discussions and ideas.}

\ccsdesc[500]{Theory of computation~Type theory}%TODO mandatory: Please choose ACM 2012 classifications from https://dl.acm.org/ccs/ccs_flat.cfm 

\keywords{Sized Types, Parametricity, Realisability, Impredicativity, Constructive Ordinals, (Co)inductive Types} %TODO mandatory; please add comma-separated list of keywords

%\category{} %optional, e.g. invited paper

%Editor-only macros:: begin (do not touch as author)%%%%%%%%%%%%%%%%%%%%%%%%%%%%%%%%%%
\EventEditors{John Q. Open and Joan R. Access}
\EventNoEds{2}
\EventLongTitle{under review for 11th FSCD (International Conference on Formal Structures for Computation and Deduction)}
\EventShortTitle{FSCD 2026}
\EventAcronym{FSCD}
\EventYear{2026}
\EventDate{July 20--23, 2026}
\EventLocation{Lisbon, Portugal}
\EventLogo{}
\SeriesVolume{}
\ArticleNo{1}
%%%%%%%%%%%%%%%%%%%%%%%%%%%%%%%%%%%%%%%%%%%%%%%%%%%%%%

\usepackage{mathtools}
\usepackage{bussproofs}
\usepackage{quiver}
\usepackage{stmaryrd}
\usepackage{wasysym}
\usepackage{capt-of}

\begin{document}
    
% General Math
\renewcommand{\.}{.\,} % Spacing
\renewcommand{\:}{:} % Spacing
\newcommand{\bN}{\mathbb{N}}
\newcommand{\isdef}{\mathrel{\overset{\makebox[0pt]{\mbox{\normalfont\tiny\sffamily def}}}{=}}}
\newcommand{\brac}[1]{\langle {#1} \rangle }
% Types
\newcommand{\Type}{~\mathsf{type}}
\newcommand{\El}{\mathrm{El}\,}
\newcommand{\Size}{\mathsf{Size}}
\newcommand{\List}{\mathsf{List} \, }
\newcommand{\SList}{\mathsf{SList} \,}
\newcommand{\nil}{\mathsf{nil} \,}
\newcommand{\cons}{\mathsf{cons}}
\newcommand{\SStream}[2]{\mathsf{Stream} \, {#1} \,[{#2}] }
\newcommand{\add}[2]{\mathsf{add}_{#1} \,x \, xs}
\newcommand{\up}{\; \uparrow \!}
\newcommand{\fix}{\mathsf{fix}}
\newcommand{\F}{F}
\newcommand{\Fmap}{\mathsf{Fmap}}
\newcommand{\map}{\mathsf{map}}
\newcommand{\fold}{\mathsf{fold}}
\newcommand{\unfold}{\mathsf{unfold}}
\newcommand{\FAlg}{F\text-\sf{Alg}}
\newcommand{\sFAlg}{F[\Diamond{-}]\text-\sf{Alg}}
\newcommand{\FCoalg}{F\text-\sf{Coalg}}
\newcommand{\sFCoalg}{F[\Box{-}]\text-\sf{CoAlg}}

\newcommand{\In}{\mathsf{in}}
\newcommand{\out}{\mathsf{out}}
\newcommand{\sto}{\overset{i}{\to}}

\newcommand{\pair}{\mathsf{pair} }
\newcommand{\inl}{\mathsf{inl}\, }
\newcommand{\inr}{\mathsf{inr} \,}
\newcommand{\Id}{\mathsf{Id}}

\newcommand{\sem}[1]{\llbracket {#1} \rrbracket}

% Category Theory
\newcommand{\cat}[1]{{\mathrm{\mathbf{#1}}}}
\newcommand{\Set}{\cat{Set}}
\newcommand{\Rel}{\cat{Rel}}
\newcommand{\Cat}{\cat{Cat}}
\newcommand{\Ptl}{\cat{Ptl}}
\newcommand{\C}{\mathcal C}
\newcommand{\E}{\mathbb E}
\newcommand{\Ob}[1]{{\mathrm{Ob}({#1})}}
\newcommand{\Hom}{\text{Hom}}
\newcommand{\id}{\mathsf{id}}

\newcommand{\fst}{\mathsf{fst}}
\newcommand{\snd}{\mathsf{snd}}

\newcommand{\ind}{\mathsf{ind}}
\newcommand{\0}{\mathbf{0}}
\newcommand{\1}{\mathbf{1}}

\newcommand{\refl}{\mathsf{refl}}

\newcommand\<\langle
\renewcommand\>\rangle
\newcommand\llb\llbracket
\newcommand\rrb\rrbracket
\newcommand\bb\mathbb
\providecommand\cal{}
\renewcommand\cal\mathcal
\renewcommand\sf\mathsf
\renewcommand\rm\mathrm
\renewcommand\frak\mathfrak

% PCA
\newcommand{\cA}{\mathcal{A}}
\renewcommand{\k}{\mathtt{k}}
\newcommand{\s}{\mathtt{s}}

% Assemblies
\newcommand{\dom}{\text{dom}}

\newcommand{\Per}{\mathsf{PER}}
\newcommand{\Perq}{\mathsf{PER_Q}}
\newcommand{\Mod}{\mathsf{Mod}}
\newcommand{\Asm}{\mathsf{Asm}}
\newcommand{\M}{\mathsf{M}}
\renewcommand{\E}{\mathsf{E}}

\newcommand{\pr}{\mathsf{pr}}
\newcommand{\Tr}[1]{{\|{#1}\|}}
\newcommand{\tr}[1]{|{#1}|}
\newcommand{\elim}{\mathsf{elim}}
\newcommand{\app}{\mathsf{app}}
\newcommand{\abs}{\mathsf{abs}}

% CwF
\newcommand{\Ty}{\mathsf{Ty}}
\newcommand{\Tm}{\mathsf{Tm}}
\newcommand{\Ctxt}{\mathsf{ctx}}
\newcommand{\p}{\mathbf{p}}
\renewcommand{\v}{\mathbf{v}}
\newcommand{\U}{U}

% Words
\newcommand{\Iff}{\hspace{6pt}\text{iff}\hspace{6pt}}
\newcommand{\If}{\hspace{6pt}\text{if}\hspace{6pt}}
\newcommand{\ie}{\hspace{6pt}\text{i.e.}\hspace{6pt}}
\renewcommand{\and}{\hspace{6pt}\text{and}\hspace{6pt}}
\newcommand{\Or}{\hspace{6pt}\text{or}\hspace{6pt}}
\newcommand{\when}{\hspace{6pt}\text{when}\hspace{6pt}}


    
    \maketitle
    
    \begin{abstract}
        To ensure decidability and consistency of its type theory, a proof assistant should only accept terminating recursive functions and productive corecursive functions.
        Most proof assistants enforce this through syntactic conditions, which can be restrictive and non-modular.
        Sized types are a type-based alternative where (co)inductive types are annotated with additional size information.
        Well-founded induction on sizes can then be used to prove termination and productivity.

        An implementation of sized types exists in \texttt{Agda}, but it is currently inconsistent due to the addition of a largest size.
        We investigate an alternative approach, where intensional type theory is extended with a large type of sizes and parametric quantifiers over sizes.
        We show that inductive and coinductive types can be constructed in this theory, which improves on earlier work where this was only possible for the finitely-branching inductive types.
        The consistency of the theory is justified by an impredicative realisability model, which interprets the type of sizes as an uncountable ordinal.
    \end{abstract}
    
    \section{Introduction}
    \label{sec:intro}
    \input{chapters/introduction}
    
    \section{The type theory} 
    \label{sec:type-theory}
    \input{chapters/theory}
    
    \section{Encoding (co)inductive types}
    \label{sec:encoding}
    \input{chapters/encoding}
    
    \section{Realisability model}
    \label{sec:model}
    \input{chapters/model}
    
    \section{Conclusion} 
    \label{sec:conclusion}
    \input{chapters/conclusion}

    \bibliography{references}
    
    \appendix
    \input{chapters/appendix}
\end{document}
